\documentclass[10pt,a4paper]{amsart}
\usepackage{amsmath,amssymb,amsthm,graphicx}
\usepackage{verbatim}
\newenvironment{reference}{\comment}{\endcomment}
%\newenvironment{reference}{}{}
\newenvironment{slogan}{\comment}{\endcomment}
\newenvironment{history}{\comment}{\endcomment}

% For commutative diagrams we use Xy-pic
\usepackage[all]{xy}

% We use 2cell for 2-commutative diagrams.
\xyoption{2cell}
% Theorem environments.
%
\theoremstyle{plain}
\newtheorem{theorem}[subsection]{Theorem}
\newtheorem{proposition}[subsection]{Proposition}
\newtheorem{lemma}[subsection]{Lemma}

\theoremstyle{definition}
\newtheorem{definition}[subsection]{Definition}
\newtheorem{example}[subsection]{Example}
\newtheorem{exercise}[subsection]{Exercise}
\newtheorem{situation}[subsection]{Situation}

\theoremstyle{remark}
\newtheorem{remark}[subsection]{Remark}
\newtheorem{remarks}[subsection]{Remarks}

\numberwithin{equation}{subsection}

% Macros
%
\def\lim{\mathop{\mathrm{lim}}\nolimits}
\def\colim{\mathop{\mathrm{colim}}\nolimits}
\def\Spec{\mathop{\mathrm{Spec}}}
\def\Hom{\mathop{\mathrm{Hom}}\nolimits}
\def\Ext{\mathop{\mathrm{Ext}}\nolimits}
\def\SheafHom{\mathop{\mathcal{H}\!\mathit{om}}\nolimits}
\def\SheafExt{\mathop{\mathcal{E}\!\mathit{xt}}\nolimits}
\def\Sch{\mathit{Sch}}
\def\Mor{\mathop{\mathrm{Mor}}\nolimits}
\def\Ob{\mathop{\mathrm{Ob}}\nolimits}
\def\Sh{\mathop{\mathit{Sh}}\nolimits}
\def\NL{\mathop{N\!L}\nolimits}
\def\CH{\mathop{\mathrm{CH}}\nolimits}
\def\proetale{{pro\text{-}\acute{e}tale}}
\def\etale{{\acute{e}tale}}
\def\QCoh{\mathit{QCoh}}
\def\Ker{\mathop{\mathrm{Ker}}}
\def\Im{\mathop{\mathrm{Im}}}
\def\Coker{\mathop{\mathrm{Coker}}}
\def\Coim{\mathop{\mathrm{Coim}}}

% Boxtimes
%
\DeclareMathSymbol{\boxtimes}{\mathbin}{AMSa}{"02}

%
% Macros for moduli stacks/spaces
%
\def\QCohstack{\mathcal{QC}\!\mathit{oh}}
\def\Cohstack{\mathcal{C}\!\mathit{oh}}
\def\Spacesstack{\mathcal{S}\!\mathit{paces}}
\def\Quotfunctor{\mathrm{Quot}}
\def\Hilbfunctor{\mathrm{Hilb}}
\def\Curvesstack{\mathcal{C}\!\mathit{urves}}
\def\Polarizedstack{\mathcal{P}\!\mathit{olarized}}
\def\Complexesstack{\mathcal{C}\!\mathit{omplexes}}
% \Pic is the operator that assigns to X its picard group, usage \Pic(X)
% \Picardstack_{X/B} denotes the Picard stack of X over B
% \Picardfunctor_{X/B} denotes the Picard functor of X over B
\def\Pic{\mathop{\mathrm{Pic}}\nolimits}
\def\Picardstack{\mathcal{P}\!\mathit{ic}}
\def\Picardfunctor{\mathrm{Pic}}
\def\Deformationcategory{\mathcal{D}\!\mathit{ef}}
\usepackage{cleveref}
\crefname{originallemma}{original lemma}{original lemmas}
\Crefname{originallemma}{Original lemma}{Original lemmas}
\makeatletter
\newlabel{lemma-characterize-separable-field-extensions}{{030W}{source}{}{sourcepointer-1}{}}
\newlabel{lemma-characterize-separable-field-extensions@cref}{{[originallemma][0][]030W}{[1][0][]source}{}{}{}}
\makeatother
\crefname{originallemma}{original lemma}{original lemmas}
\Crefname{originallemma}{Original lemma}{Original lemmas}
\makeatletter
\newlabel{lemma-generic-points-geometrically-reduced}{{07K2}{source}{}{sourcepointer-2}{}}
\newlabel{lemma-generic-points-geometrically-reduced@cref}{{[originallemma][0][]07K2}{[1][0][]source}{}{}{}}
\makeatother
\crefname{originallemma}{original lemma}{original lemmas}
\Crefname{originallemma}{Original lemma}{Original lemmas}
\makeatletter
\newlabel{lemma-limit-argument}{{00I3}{source}{}{sourcepointer-3}{}}
\newlabel{lemma-limit-argument@cref}{{[originallemma][0][]00I3}{[1][0][]source}{}{}{}}
\makeatother
\crefname{originallemma}{original lemma}{original lemmas}
\Crefname{originallemma}{Original lemma}{Original lemmas}
\makeatletter
\newlabel{lemma-make-separably-generated}{{04KM}{source}{}{sourcepointer-4}{}}
\newlabel{lemma-make-separably-generated@cref}{{[originallemma][0][]04KM}{[1][0][]source}{}{}{}}
\makeatother
\crefname{originallemma}{original lemma}{original lemmas}
\Crefname{originallemma}{Original lemma}{Original lemmas}
\makeatletter
\newlabel{lemma-minimal-prime-reduced-ring}{{00EU}{source}{}{sourcepointer-5}{}}
\newlabel{lemma-minimal-prime-reduced-ring@cref}{{[originallemma][0][]00EU}{[1][0][]source}{}{}{}}
\makeatother
\crefname{originallemma}{original lemma}{original lemmas}
\Crefname{originallemma}{Original lemma}{Original lemmas}
\makeatletter
\newlabel{lemma-separable-extension-preserves-reducedness}{{030U}{source}{}{sourcepointer-6}{}}
\newlabel{lemma-separable-extension-preserves-reducedness@cref}{{[originallemma][0][]030U}{[1][0][]source}{}{}{}}
\makeatother
\title{Five equivalent tests for geometric reducedness via purely inseparable extensions and closures}
\author{The Stacks Project Authors}
\date{stacks-project:algebra.tex:local-snapshot-20261012; GNU FDL 1.2 or later}
\begin{document}
\maketitle
\noindent\textit{Editorial note.} All delivered mathematical text is the author's own, copied from the pinned Stacks Project chapter file algebra.tex; no statement or proof step was rewritten, reordered or removed.\par\medskip
\noindent\textit{Editorial note.} Preamble provenance, as required for this source: the frozen snapshot stages the chapter's own preamble file beside it (preamble.tex, 217 lines, 8116 bytes, the same shared Stacks-wide preamble the published chow and etale-cohomology cards use); the first line of the pinned chapter file is \textbackslash{}input\{preamble\}, so no preamble text lives inside the chapter itself. The card takes that file's lines 8-18 (the verbatim package and the three \textbackslash{}comment-based environments reference/slogan/history at lines 9-12, section 13.2, together with the Xy-pic setup) and its lines 154 to the runtime end of the file (the native theorem declarations and mathematical macros). Not one character was changed, no package was added, no macro was defined, and the excerpt serves only as preamble, never as a span coordinate.\par\medskip
\noindent\textit{Editorial note.} The same-chapter pointers are located at their \textbackslash{}label lines inside the pinned chapter file; the cross-chapter pointers are located in the project's own tag list staged in the snapshot (sources/stacks-chow/tags.txt), at the verbatim lines carrying their canonical tags.\par\medskip
\noindent\textit{Editorial note.} The counted proof refers to same-chapter statements by label; they are recorded as navigation stubs pointed at their label lines.\par\medskip
\noindent\textit{Editorial note.} The counted spans contain no cite, so the card carries no bibliography.\par\medskip
\noindent\textit{Allowed standard prerequisites (editorial).}\begin{itemize}
\item Embedding finite purely inseparable extensions in the perfect closure. Used at the start of the counted proof (lines 10658-10662). Orders the five conditions so that the implications in one direction are immediate.
\item Separable field extensions tested by a characterisation. Invoked near the end of the counted proof (lines 10692-10694). Supplies the separability needed to conclude reducedness of the base change.
\end{itemize}\medskip
\ifdefined\hypertarget\hypertarget{sourcepointer-1}{}\fi
\noindent\textit{Source reference.} Original lemma 030W: sources/stacks-algebra/pinned-original-source.tex, source lines 10577--10577. Navigation pointer for \textbackslash{}ref\{lemma-characterize-separable-field-extensions\} used inside the counted span (line 10692); the argument that consumes it is retained in the span.\par\medskip
\ifdefined\hypertarget\hypertarget{sourcepointer-2}{}\fi
\noindent\textit{Source reference.} Original lemma 07K2: sources/stacks-algebra/pinned-original-source.tex, source lines 10414--10414. Navigation pointer for \textbackslash{}ref\{lemma-generic-points-geometrically-reduced\} used inside the counted span (line 10693); the argument that consumes it is retained in the span.\par\medskip
\ifdefined\hypertarget\hypertarget{sourcepointer-3}{}\fi
\noindent\textit{Source reference.} Original lemma 00I3: sources/stacks-algebra/pinned-original-source.tex, source lines 10323--10323. Navigation pointer for \textbackslash{}ref\{lemma-limit-argument\} used inside the counted span (line 10669); the argument that consumes it is retained in the span.\par\medskip
\ifdefined\hypertarget\hypertarget{sourcepointer-4}{}\fi
\noindent\textit{Source reference.} Original lemma 04KM: sources/stacks-algebra/pinned-original-source.tex, source lines 10179--10179. Navigation pointer for \textbackslash{}ref\{lemma-make-separably-generated\} used inside the counted span (line 10677); the argument that consumes it is retained in the span.\par\medskip
\ifdefined\hypertarget\hypertarget{sourcepointer-5}{}\fi
\noindent\textit{Source reference.} Original lemma 00EU: sources/stacks-algebra/pinned-original-source.tex, source lines 4760--4760. Navigation pointer for \textbackslash{}ref\{lemma-minimal-prime-reduced-ring\} used inside the counted span (line 10690); the argument that consumes it is retained in the span.\par\medskip
\ifdefined\hypertarget\hypertarget{sourcepointer-6}{}\fi
\noindent\textit{Source reference.} Original lemma 030U: sources/stacks-algebra/pinned-original-source.tex, source lines 10369--10369. Navigation pointer for \textbackslash{}ref\{lemma-separable-extension-preserves-reducedness\} used inside the counted span (line 10679); the argument that consumes it is retained in the span.\par\medskip
\begin{lemma}
\label{lemma-geometrically-reduced-finite-purely-inseparable-extension}
Let $k$ be a field. Let $S$ be a $k$-algebra.
The following are equivalent:
\begin{enumerate}
\item $k' \otimes_k S$ is reduced for every finite
purely inseparable extension $k'$ of $k$,
\item $k^{1/p} \otimes_k S$ is reduced,
\item $k^{perf} \otimes_k S$ is reduced, where $k^{perf}$ is the
perfect closure of $k$,
\item $\overline{k} \otimes_k S$ is reduced, where $\overline{k}$ is the
algebraic closure of $k$, and
\item $S$ is geometrically reduced over $k$.
\end{enumerate}
\end{lemma}

\begin{proof}
Note that any finite purely inseparable extension $k'/k$ embeds
in $k^{perf}$. Moreover, $k^{1/p}$ embeds into $k^{perf}$ which embeds
into $\overline{k}$. Thus it is
clear that (5) $\Rightarrow$ (4) $\Rightarrow$ (3) $\Rightarrow$ (2)
and that (3) $\Rightarrow$ (1).

\medskip\noindent
We prove that (1) $\Rightarrow$ (5).
Assume $k' \otimes_k S$ is reduced for every finite
purely inseparable extension $k'$ of $k$. Let $K/k$ be
an extension of fields. We have to show that $K \otimes_k S$
is reduced. By Lemma \ref{lemma-limit-argument} we reduce to the case where
$K/k$ is a finitely generated field extension. Choose a diagram
$$
\xymatrix{
K \ar[r] & K' \\
k \ar[u] \ar[r] & k' \ar[u]
}
$$
as in Lemma \ref{lemma-make-separably-generated}.
By assumption $k' \otimes_k S$ is reduced.
By Lemma \ref{lemma-separable-extension-preserves-reducedness}
it follows that $K' \otimes_k S$ is reduced.
Hence we conclude that $K \otimes_k S$ is reduced as desired.

\medskip\noindent
Finally we prove that (2) $\Rightarrow$ (5).
Assume $k^{1/p} \otimes_k S$ is reduced. Then $S$ is reduced.
Moreover, for each localization $S_{\mathfrak p}$ at a minimal
prime $\mathfrak p$, the ring $k^{1/p}\otimes_k S_{\mathfrak p}$
is a localization of $k^{1/p} \otimes_k S$ hence is reduced.
But $S_{\mathfrak p}$ is a field by
Lemma \ref{lemma-minimal-prime-reduced-ring},
hence $S_{\mathfrak p}$ is geometrically reduced by
Lemma \ref{lemma-characterize-separable-field-extensions}.
It follows from Lemma \ref{lemma-generic-points-geometrically-reduced}
that $S$ is geometrically reduced.
\end{proof}

\end{document}
